\section{The original canonical tower as an elementary principalization}

The compact ambient manifold in this section may be disconnected, and may
be empty.  The source of a smooth base map need not be compact.  We use the
literal marked ideal $M=(X,X,\varnothing,I,1)$ throughout.  The finite tower
is the actual output of the compact canonical driver for this record.
Thus the elementary principalization and the original terminal manifold
refer to the same resolution, including its full total-transform ideal.
The construction specializes the terminal-map comparison to elementary
principalization, as in \cite{bierstone-milman08-functoriality}; it does not
identify an embedded-resolution stopped prefix with a full terminal tower.

\begin{theorem}[Canonical elementary principalization on the actual compact terminal tower]
\label{hir:original-canonical-elementary-principalization}
\uses{hir:principalization,hir:principalization-compact-support,hir:admissible-sequence,hir:smooth-original-compact-canonical-terminal-object,hir:original-global-terminal-finite-whole-biholomorphism}
\lean{Hironaka.HironakaOriginalCanonicalElementaryPrincipalization.canonical_elementary_principalization}
\leanok
\source{bierstone-milman08-functoriality,wlodarczyk05-hironaka}
Let $X$ be a compact complex manifold and $I$ a coherent ideal which is
nonzero as a germ at every point.  Let $T:D\to X$ be the internally
constructed compact canonical tower for $M$.  Write $p=T.\mathrm{proj}$,
$D_T$ for its effective exceptional divisor and $E_T$ for the actual
accumulated boundary of the transform of $M$.  Then $T$ is both a canonical
segment and a geometric resolution of $M$.  The same projection $p$ is a
proper surjective holomorphic admissible blowup sequence from the empty
boundary, every centre lies over $Z(I)$, and
\[
 p^* I=I_{D_T},\qquad E_T\text{ is SNC},\qquad
 \operatorname{supp}D_T=\bigcup_{H\in E_T}H=p^{-1}Z(I).
\]
Its exceptional locus is contained in $\operatorname{supp}D_T$, and it is
an isomorphism over $X\setminus Z(I)$.  There is a biholomorphism
$\eta_X:\widetilde X_M\to D$ from the original glued terminal manifold
such that $p\eta_X=P_M$, the original whole glued projection.
\end{theorem}
\begin{proof}
Coherence and germwise nonvanishing put the literal $M$ in the ambient
invariant domain.  Apply the compact canonical driver to this very record;
retain its canonical segment and geometric resolution.  The segment gives
whole canonicality.  The original cosupport is exactly $Z(I)$, so the
actual test-tower centre containment puts every centre of this tower over
$Z(I)$.  Applying the admissible-sequence dictionary to this tower gives
the claimed sequence with exactly its accumulated boundary.  Its
projection is holomorphic, proper and surjective because it is the
projection of a finite actual blowup tower.

For mark one and ambient support $X$, the geometric-resolution identity
gives $p^*I=I_{D_T}$ as whole germ ideals.  Well-formedness of the actual
transformed record gives SNC for $E_T$.  The tower's exceptional-divisor
support formula, with empty initial boundary, gives the first support
equality.  The zero-locus pullback formula applied to the whole ideal
identity gives the second.  The exceptional-locus containment for a tower
therefore puts its entire image over $Z(I)$; the blowup-sequence
isomorphism theorem gives the isomorphism over its complement.  Finally
apply the original finite-terminal identification to this same $T$ and
the canonicality and resolution just produced.  Its whole projection
formula is precisely $p\eta_X=P_M$.
\end{proof}

\begin{theorem}[Arbitrary smooth source and the same elementary principalization]
\label{hir:original-canonical-elementary-principalization-smooth}
\uses{hir:original-canonical-elementary-principalization,hir:original-global-finite-smooth-source-history,hir:smooth-original-finite-terminal-intrinsic-smooth-map,hir:original-global-finite-terminal-smooth-map-cartesian,hir:original-global-terminal-finite-whole-biholomorphism}
\lean{Hironaka.HironakaOriginalCanonicalElementaryPrincipalization.canonical_elementary_principalization_smooth_base_change}
\leanok
\source{bierstone-milman08-functoriality}
With the preceding original $I,X,M,T,p,D_T$, let $b:Y\to X$ be any
holomorphic submersion.  The actual Cartesian pullback construction of
this finite $T$ internally gives a finite canonical geometric resolution
$S:C\to Y$ of $b^*M$, without compactness of $Y$.  Writing $q=S.\mathrm{proj}$,
there are internally constructed holomorphic submersions
$g:C\to D$ and $G:\widetilde Y_{b^*M}\to\widetilde X_M$ for which
\[
 pg=bq,\qquad P_MG=bP_{b^*M},\qquad
 \eta_XG=g\eta_Y .
\]
Both displayed projection squares are holomorphic Cartesian squares;
$\eta_Y:\widetilde Y_{b^*M}\to C$ and $\eta_X$ are biholomorphisms over
the unchanged original bases.  In particular the lower principalization
in these squares is exactly the preceding $p$, with its literal whole
identity $p^*I=I_{D_T}$.
\end{theorem}
\begin{proof}
Choose the actual smooth pullback tower of the entire finite target $T$.
The original smooth-source-history theorem proves canonicality of this
actual $S$ and its genuine geometric terminality.  Apply the intrinsic
finite-terminal smooth-map theorem to these two histories, keeping the
already chosen target $T$.  It gives the submersion $g$, the equality
$pg=bq$, and the holomorphic Cartesian universal property.

Identify both actual finite terminal histories with their original glued
terminal manifolds by their genuine over-base biholomorphisms $\eta_Y$
and $\eta_X$.  The global finite smooth arrow is defined as
$G=\eta_X^{-1}g\eta_Y$.  Hence $\eta_XG=g\eta_Y$.  The finite comparison
projection identities give the second projection equality.  Transporting
the finite Cartesian universal property through these biholomorphisms
gives the whole global holomorphic Cartesian property, and their
composition with $g$ is a submersion.  The target is still the original
compact driver's tower, so the first theorem supplies its unchanged whole
total-transform identity.
\end{proof}

\begin{remark}
These results give the arbitrary smooth comparison for elementary ideal
principalization with empty initial boundary.  Embedded resolution uses a
stopped prefix before the first centre meeting the original regular locus;
the full terminal ideal map here is not that stopped map.  Comparisons for
that prefix and the corresponding boundary statements remain separate.
\end{remark}
